\section{Stakes: the lineage panel, scoped}
\label{sec:stakes}

% source: PAPER_SKELETON Sec 8; EVIDENCE_lineage_boundaries L-lab / tier table.

\paragraph{The lineage panel is descriptive and scoped, not a vendor
leaderboard.} The cross-lineage comparison of Section~\ref{sec:production}
uses one seed, is single-format only where it compares directly, and is
anchored through the format-robust echo-corrected reuse metric. We present
it as stakes, why in-stream absorption matters for chain-of-thought
monitoring, with the universal recompute column as its non-partisan center,
rather than as a ranking.
% source: PAPER_SKELETON 8.1; EVIDENCE_lineage_boundaries L-lab.1 caveat

\paragraph{The load-bearing comparison is within one format.} The single
most defensible statement is the within-format one: same instruct format,
each model's own traces, audited-false plants, gpt-4.1 reuses a false
category at 0.58 while claude-sonnet-5 reuses at 0.01.
% source: PAPER_SKELETON 8.2; EVIDENCE_lineage_boundaries L-lab.1
We never say the Anthropic models check every time; the supported claim is
that they reject every audited falsehood in this harness, at closure-valid
0.97 to 1.00 and reuse 0.000, and the verbatim one-hop-check transcript is
attributed to a single model. We never say gpt-4.1 never checks; only its
reuse rate is on record.
% source: EVIDENCE_lineage_boundaries L-lab.2 caveat; STYLE_SHEET Sec 5

\paragraph{The split is not memorization.} The worlds use nonce predicates,
and the same Anthropic models that reject re-readable plants absorb
recomputable ones at 1.00, which a memorization account cannot produce.
% source: PAPER_SKELETON 8.3; EVIDENCE_lineage_boundaries L-rec.1 / L-lab.2

\paragraph{Every by-model claim carries an evidence tier.} To keep the
descriptive arm from being read as a controlled result, Table~\ref{tab:tiers}
labels each claim by how much weight it carries: controlled on the anchor
model, controlled multi-model replication, descriptive single-seed API, or
legacy appendix. The replication tier now carries the non-replication
verdict, and the anchor-model tier is explicitly scoped to Qwen2.5-7B where
the replication tested it.
% source: PAPER_SKELETON 8.4; EVIDENCE_lineage_boundaries tier table + E1 landing note

\begin{table}[t]
\centering
\caption{\textbf{Evidence-tier map.} Each family of claims is labeled by how
much weight it carries. Tier~2 now holds a verified non-replication, and
tier~1 is scoped to Qwen2.5-7B where tier~2 tested it. Cumulative API spend
across the descriptive arm was \$30.81.}
% source: EVIDENCE_lineage_boundaries tier table
\label{tab:tiers}
\small
\begin{tabular}{@{}llp{7.2cm}@{}}
\toprule
Tier & Basis & Claims \\
\midrule
1 anchor (controlled) & Qwen2.5-7B, predictions fixed in advance & Visibility cliff, usability asymmetry, lexical cueing. Scoped to Qwen2.5-7B per tier~2. \\
2 replication (controlled) & multi-model, predictions fixed in advance & \textbf{Non-replication (verified): 0/4 fully replicate; visibility cliff reverses on 3/4; usability gap holds on 2/4 but collapses at 32B; flatness 0/4.} \\
3 descriptive API & single-seed & Production boundary (corroboration), lab split, recompute boundary. \\
4 legacy / appendix & instrument lesson & Certificate experiment, legacy gap decomposition. \\
\bottomrule
\end{tabular}
\end{table}
